\documentclass[10pt,letterpaper]{article}
\usepackage{cornell-notes}

\title{Clause 27.08.04.02: Lower Bound}
\author{}
\date{}
\setDocTitle{Clause 27.08.04.02: Lower Bound}
\setDocSubtitle{C++ 2024}
\setDocOwner{C++ 2024 Cornell Notes}
\setDocDate{\today}
\setCornellCollection{C++ 2024 Cornell Notes}
\setCornellUnitType{Clause}
\setCornellUnitNumber{27.08.04.02}
\setCornellUnitTitle{Lower Bound}
\hypersetup{pdftitle={Clause 27.08.04.02: Lower Bound},pdfsubject={Cornell notes},pdfauthor={}}

\begin{document}
\maketitle
\begin{CornellOverviewBox}{Overview}
\textbf{Classification:} Standard library - Normative\\[2pt]
\textbf{Learning objective:} Understand the rules, terminology, and semantic consequences associated with lower\_bound.
\end{CornellOverviewBox}\begin{CornellSummaryBox}{Summary}
{\sffamily\bfseries\color{Accent} Summary}\par\vspace{3pt}Section 27.8.4.2 focuses on lower\_bound. Apply the specified interface together with its mandates, effects, result conditions, exception behavior, and complexity constraints. When applying it, preserve the distinction between mandatory requirements, recommendations, permissions, and behavior deliberately left undefined, unspecified, or implementation-defined.
\end{CornellSummaryBox}

\end{document}
